% ==============================================================================
\section{MILP Formulation}
\label{sec:milp}
% ==============================================================================

We present a mixed-integer linear programming formulation of the
problem of \S\ref{sec:problem}.  Where possible the notation follows
the conventions used elsewhere in the literature for aircraft
positioning problems with blocking arcs: parameters use upper-case
Latin letters, decision variables use lower-case Latin letters, and
Greek letters are reserved for fixed scalar constants ($\varepsilon$,
$\mu$, $\delta$, the time-granularity $\eta$ and the big-$M$
relaxation~$M$).  The key feature of the formulation is that every
access of a blocked aircraft is explicitly \emph{classified}, by means
of a partition constraint, into one of the four mutually exclusive
states allowed by the problem: before the front aircraft's stay,
after it, in an inter-job gap (Mode~B), or strictly inside an
interruptible job (Mode~C).  An access that cannot be classified into
any of these four states corresponds to an attempt to enter or leave a
position whose blocker is in the middle of a non-interruptible job,
and is correctly rendered infeasible by the model.

\subsection*{Sets}

\begin{tabular}{@{}l@{\quad}p{0.78\textwidth}@{}}
  $\calP$
    & maintenance positions \\
  $\calR$
    & aircraft \\
  $\calJ_r,\ r \in \calR$
    & ordered job chain of aircraft $r$;
      $\calJ_r = (j^r_1, \ldots, j^r_{N_r})$ \\
  $\calJ = \bigcup_{r \in \calR} \calJ_r$
    & global job set \\
  $\calA \subseteq \calP \times \calP$
    & blocking arcs $(p, p')$: position $p$ blocks position $p'$ \\
\end{tabular}

\subsection*{Parameters}

\begin{tabular}{@{}l@{\quad}p{0.78\textwidth}@{}}
  $D_j$
    & nominal duration of job $j \in \calJ$ \\
  $I_j \in \{0, 1\}$
    & interruptibility flag of job $j$; $I_j = 1$ iff $j$ can be paused \\
  $E_r, L_r$
    & earliest start and target finish of aircraft $r$ \\
  $\varepsilon$
    & tow-in / tow-out time; minimum separation between consecutive
      aircraft serviced at the same position \\
  $\mu$
    & operational time absorbed by a single Mode-(B) manoeuvre inside
      an inter-job gap (per-access, cumulative across all manoeuvres
      that share the same gap) \\
  $\delta$
    & additional working time added to an interruptible job each time
      it is paused to allow a Mode-(C) manoeuvre \\
  $\eta$
    & time-granularity used to enforce strict inequalities on the
      Mode-(C) job interior; equals the smallest time difference
      meaningful in the instance (e.g.\ one day) \\
  $K^{\max}_j$
    & valid upper bound on the number of times job $j$ may be paused;
      we take $K^{\max}_j = 2 (|\calR| - 1)$ \\
  $H^{UB}$
    & technical upper bound on every job start and finish time,
      large enough to accommodate the worst-case sequential schedule
      with all possible interruptions \\
  $M$
    & big-$M$ constant; we take $M = H^{UB} + \max(\varepsilon, \mu, \delta, \eta)$ \\
  $W^M, W^D, W^S$
    & weights for makespan, total delay, and movements \\
\end{tabular}

\subsection*{Variables}

\paragraph{Assignment and timing.}
Binaries $x_{rp} \in \{0,1\}$ (aircraft $r$ assigned to position $p$)
and $q_{rr'p} \in \{0,1\}$ (aircraft $r$ precedes aircraft $r'$ at
$p$).  Continuous $s_j, f_j \in [0, H^{UB}]$ for every $j \in \calJ$,
together with the aircraft-level shortcuts $s_r = s_{j^r_1}$ and
$f_r = f_{j^r_{N_r}}$.  The integer counter $\kappa_j \in \{0, \ldots,
K^{\max}_j\}$ records how many times job $j$ is paused.

\paragraph{Placement compatibility.}
For every $r \ne r'$ and every $(p, p') \in \calA$ we introduce a
binary $y_{rr'pp'} \in \{0,1\}$ that equals $1$ iff $r$ is the front
aircraft at $p$ and $r'$ is the rear aircraft at $p'$.

\paragraph{Access classification.}
For every $r \ne r'$, every $(p, p') \in \calA$, and every access
side $a \in \{\mathrm{in}, \mathrm{out}\}$ (with the convention $\tau^a_{r'}
= s_{r'}$ when $a = \mathrm{in}$ and $\tau^a_{r'} = f_{r'}$ when $a =
\mathrm{out}$), we introduce the four mutually exclusive classification
indicators
\begin{align*}
  & z^{-,a}_{rr'pp'} \in \{0,1\}
    && \text{access $\tau^a_{r'}$ falls strictly before $r$'s stay,} \\
  & z^{+,a}_{rr'pp'} \in \{0,1\}
    && \text{access $\tau^a_{r'}$ falls strictly after $r$'s stay,} \\
  & b^{Ba}_{rr'pp'k} \in \{0,1\},\ k = 1, \ldots, N_r-1
    && \text{access falls in inter-job gap $k$ of $r$ (Mode B),} \\
  & b^{Ca}_{rr'pp'j} \in \{0,1\},\ j \in \calJ_r
    && \text{access falls strictly inside job $j$ of $r$ (Mode C).}
\end{align*}

\paragraph{Objective auxiliaries.}
Makespan $m \in \NR$, per-aircraft delay $v^D_r \in \NR$, and total
movement count $n \in \NZ$.

\subsection*{Constraints}

To keep the alignment compact we abbreviate
$\forall r \ne r' \in \calR,\ (p,p') \in \calA,\ a \in
\{\mathrm{in},\mathrm{out}\}$ by $(\dagger)$ and
$\forall r \ne r' \in \calR,\ p \in \calP$ by $(\ddagger)$.

\paragraph{Assignment, chain, earliest start, sequencing.}
\begin{align}
  \sum_{p \in \calP} x_{rp} &= 1
    && \forall r \in \calR \label{eq:assign}\\
  f_j &= s_j + D_j + \delta\, \kappa_j
    && \forall j \in \calJ \label{eq:jobdur}\\
  s_{j^r_{k+1}} &\ge f_{j^r_k}
    && \forall r,\ k = 1, \ldots, N_r{-}1 \label{eq:chain}\\
  s_r &\ge E_r
    && \forall r \in \calR \label{eq:earlystart}\\
  s_{r'} &\ge f_r + \varepsilon - M(1 - q_{rr'p}) \notag\\
         &\quad - M(1 - x_{rp}) - M(1 - x_{r'p})
    && (\ddagger) \label{eq:sep}\\
  q_{rr'p} + q_{r'rp} &\ge x_{rp} + x_{r'p} - 1
    && \forall r < r' \in \calR,\, p \in \calP \label{eq:ordLB}\\
  q_{rr'p} + q_{r'rp} &\le 1
    && \forall r < r' \in \calR,\, p \in \calP \label{eq:ordUB}\\
  q_{rr'p} &\le x_{rp},\quad q_{rr'p} \le x_{r'p}
    && (\ddagger),\ \forall p \in \calP \label{eq:ordTight}
\end{align}

\paragraph{Placement compatibility.}
For every $r \ne r' \in \calR$ and every $(p, p') \in \calA$:
\begin{align}
  y_{rr'pp'} &\le x_{rp},\quad
  y_{rr'pp'} \le x_{r'p'},\quad
  y_{rr'pp'} \ge x_{rp} + x_{r'p'} - 1
    \label{eq:y_link}
\end{align}

\paragraph{Exhaustive classification of every access.}
For every $(\dagger)$, an access $\tau^a_{r'}$ that is exposed to the
front aircraft $r$ must be classified into exactly one of the four
allowed states; conversely, if $y_{rr'pp'} = 0$ the access is not
exposed and no indicator may fire:
\begin{align}
  z^{-,a}_{rr'pp'} + z^{+,a}_{rr'pp'}
   + \!\!\sum_{k=1}^{N_r-1}\! b^{Ba}_{rr'pp'k}
   + \!\!\sum_{j \in \calJ_r}\! b^{Ca}_{rr'pp'j}
   &= y_{rr'pp'}
    && \label{eq:partition}
\end{align}

\paragraph{Before / after the front stay.}
For every $(\dagger)$:
\begin{align}
  \tau^a_{r'} &\le s_r - \eta + M\bigl(1 - z^{-,a}_{rr'pp'}\bigr)
    && \label{eq:zminus}\\
  \tau^a_{r'} &\ge f_r + \eta - M\bigl(1 - z^{+,a}_{rr'pp'}\bigr)
    && \label{eq:zplus}
\end{align}

\paragraph{Mode (B) --- inter-job gap.}
For every $(\dagger)$ and every $k = 1, \ldots, N_r-1$:
\begin{align}
  \tau^a_{r'} &\ge f_{j^r_k} - M\bigl(1 - b^{Ba}_{rr'pp'k}\bigr)
    && \label{eq:bb_lb}\\
  \tau^a_{r'} &\le s_{j^r_{k+1}} + M\bigl(1 - b^{Ba}_{rr'pp'k}\bigr)
    && \label{eq:bb_ub}
\end{align}
Each Mode-(B) event consumes $\mu$ time units of the gap, so the gap
must be wide enough to accommodate all events that use it:
\begin{align}
  s_{j^r_{k+1}} - f_{j^r_k}
    &\ge
    \mu \!\!\sum_{\substack{r' \ne r\\ (p,p') \in \calA \\ a \in \{\mathrm{in}, \mathrm{out}\}}}
        \!\! b^{Ba}_{rr'pp'k}
    && \forall r \in \calR,\ k = 1, \ldots, N_r-1
    \label{eq:bb_gap}
\end{align}

\paragraph{Mode (C) --- strictly inside an interruptible job.}
For every $(\dagger)$ and every $j \in \calJ_r$:
\begin{align}
  \tau^a_{r'} &\ge s_j + \eta - M\bigl(1 - b^{Ca}_{rr'pp'j}\bigr)
    && \label{eq:bc_lb}\\
  \tau^a_{r'} &\le s_j + D_j + \delta\bigl(\kappa_j - b^{Ca}_{rr'pp'j}\bigr)
                - \eta + M\bigl(1 - b^{Ca}_{rr'pp'j}\bigr)
    && \label{eq:bc_ub}\\
  b^{Ca}_{rr'pp'j} &\le I_j
    && \label{eq:bc_interr}
\end{align}
The upper bound~\eqref{eq:bc_ub} pins $\tau^a_{r'}$ inside the
\emph{active} portion of the job, excluding the extension that this
specific access would itself generate, and so prevents the
``self-justifying interruption'' phenomenon in which an access could
otherwise become feasible only because the very interruption it
triggers extends the job past it.  Constraint~\eqref{eq:bc_interr}
forbids Mode~(C) on non-interruptible jobs.

\paragraph{Counters, delay and makespan.}
\begin{align}
  \kappa_j &=
    \!\!\sum_{\substack{r' \ne r\\ (p,p') \in \calA\\ a \in \{\mathrm{in}, \mathrm{out}\}}}
    \!\! b^{Ca}_{rr'pp'j}
    && \forall r \in \calR,\ j \in \calJ_r \label{eq:kappa}\\
  n &= 2 \!\!\sum_{(p,p') \in \calA}
       \!\!\sum_{\substack{r,r' \in \calR \\ r \ne r'}}
       \sum_{a}
       \biggl(
         \!\sum_{k=1}^{N_r-1}\! b^{Ba}_{rr'pp'k}
         + \!\sum_{j \in \calJ_r}\! b^{Ca}_{rr'pp'j}
       \biggr)
    && \label{eq:movements}\\
  v^D_r &\ge f_r - L_r
    && \forall r \in \calR \label{eq:delay}\\
  m &\ge f_r
    && \forall r \in \calR \label{eq:makespan}
\end{align}

\subsection*{Objective}
\begin{equation}
  \min \quad W^M\, m
         \;+\; W^D \sum_{r \in \calR} v^D_r
         \;+\; W^S\, n.
  \label{eq:milpobj}
\end{equation}

\paragraph{Comments on the formulation.}
Constraints~\eqref{eq:assign}--\eqref{eq:ordTight} reproduce the
assignment, chain, earliest-start and same-position sequencing
requirements: \eqref{eq:jobdur} embeds the per-job extension $\delta
\kappa_j$ into the duration; \eqref{eq:chain} together with
$s_r = s_{j^r_1}$, $f_r = f_{j^r_{N_r}}$ chains consecutive jobs of
the same aircraft; \eqref{eq:sep}--\eqref{eq:ordTight} enforce the
disjunction between two aircraft sharing a position and tighten the
ordering binaries with the standard $q \le x$ reinforcements.
Inequalities~\eqref{eq:y_link} make $y_{rr'pp'}$ the logical
conjunction of $x_{rp}$ and $x_{r'p'}$, so $y_{rr'pp'} = 1$ exactly
when the blocker / blocked placement of the arc $(p, p')$ is
materialised by the assignment.

Equation~\eqref{eq:partition} is the central structural constraint:
whenever the placement is materialised, every access $\tau^a_{r'}$
must be classified into exactly one of the four allowed states; when
the placement is not materialised, no indicator may fire.  This is
what turns the formulation into an exact model: an access that falls
inside a non-interruptible job cannot be put in Mode~(C) (forbidden
by \eqref{eq:bc_interr}), cannot be put in Mode~(B) (would violate
\eqref{eq:bb_lb}--\eqref{eq:bb_ub}), and cannot be put in $z^-$ or
$z^+$ (would violate \eqref{eq:zminus}--\eqref{eq:zplus}), so the only
way to satisfy~\eqref{eq:partition} is to revise the timing or the
assignment.  Constraint~\eqref{eq:bb_gap} forces the gap to accommodate
the cumulative manoeuvre time of all Mode-(B) events that use it, and
\eqref{eq:bc_ub} guards against self-justifying interruptions.
Equation~\eqref{eq:kappa} aggregates the Mode-(C) indicators into the
per-job counter used by~\eqref{eq:jobdur}, and~\eqref{eq:movements}
adds up all manoeuvres with the displacement-plus-return factor of
two.  The objective~\eqref{eq:milpobj} combines makespan, total
positive delay past the target finish, and movement count under
non-negative user weights; the choice of weight profile is discussed
in \S\ref{sec:exp_setup}.
